% --- Archon physics formalization source begin ---
% archon:physics
% archon:covers PhyXMiniProblems/problem_phyx_mini_0704.lean
% archon:source-report reports/phyx_mini/problem_phyx_mini_0704.source.json

\chapter{Physics problem phyx\_mini\_0704}
\label{ch:PhyXMiniProblems_problem_phyx_mini_0704}

\paragraph{Problem source.}
\#\# Physical scenario

The engine in a small airplane is specified to have a torque of \textbackslash{}(60\textbackslash{},\textbackslash{}mathrm\{N\textbackslash{},m\}\textbackslash{}). This engine drives a \textbackslash{}(2.0\textbackslash{},\textbackslash{}mathrm\{m\}\textbackslash{})-long, \textbackslash{}(40\textbackslash{},\textbackslash{}mathrm\{kg\}\textbackslash{}) propeller.

\#\# Question

On start-up,how long does it take the propeller to reach \textbackslash{}(200\textbackslash{},\textbackslash{}mathrm\{rpm\}\textbackslash{})?

\#\# Answer choices

- A: A: 5.0s

- B: B: 4.8s

- C: C: 4.6s

- D: D: 4.4s

\#\# Figure caption (auxiliary; use the image as primary evidence)

The image illustrates a propeller with an axis of rotation. The propeller is vertically oriented, and the diagram highlights various details about it:

- The propeller has a length labeled as \textbackslash{}(L = 2.0 \textbackslash{}, \textbackslash{}text\{m\}\textbackslash{}).
- It has a mass denoted as \textbackslash{}(M = 40 \textbackslash{}, \textbackslash{}text\{kg\}\textbackslash{}).
- An arrow indicates the direction of rotation due to torque from the engine.
- The phrase "The torque from the engine rotates the propeller" is written in blue text with a dotted arrow pointing towards the propeller to indicate the application of torque.
- The rotation axis is marked as "Axis".

Overall, the diagram is a clear representation of the mechanical aspects involved in the rotation of a propeller.

\#\# Dataset metadata

Rotational Motion; Physical Model Grounding Reasoning, Multi-Formula Reasoning

\paragraph{Recorded answer/context.}
C: C: 4.6s

\paragraph{Figure/image path.}
/root/proposal\_for\_physic/science-mango/phyx\_mini\_run/phyx\_data/test\_image/704.png

\paragraph{Formalization target.}
create a compiling Lean file with sorry bodies at `PhyXMiniProblems/problem\_phyx\_mini\_0704.lean`. The Lean declarations must preserve the physical quantities, dimensions or dimensional roles, figure labels, governing-law hypotheses, and final relation expressed by this problem.
Use Mathlib/Physlib names found through LeanExplore where available. If a physics API is missing, introduce faithful local abstractions rather than scalar placeholder aliases.

\begin{theorem}[Physics formalization target]
\label{thm:physics:phyx_mini_0704:target}
\lean{PhyXMiniProblems.ProblemPhyXMini0704.propeller_startup_duration_and_answer_choice}
\uses{def:physics:phyx-mini-0704:phyxminiproblems-problemphyxmini0704-timeinseconds, def:physics:phyx-mini-0704:phyxminiproblems-problemphyxmini0704-propellerstartupsetup, def:physics:phyx-mini-0704:phyxminiproblems-problemphyxmini0704-matchesproblemdata, def:physics:phyx-mini-0704:phyxminiproblems-problemphyxmini0704-matchessourcefigure, def:physics:phyx-mini-0704:phyxminiproblems-problemphyxmini0704-validpropellerstartupphysics, def:physics:phyx-mini-0704:phyxminiproblems-problemphyxmini0704-answerchoice, def:physics:phyx-mini-0704:phyxminiproblems-problemphyxmini0704-answerchoice-durationseconds}
The assigned autoformalize agent should translate this physics problem into Lean declarations in the covered file, with theorem and lemma proof bodies written as `by sorry`.
\end{theorem}
\begin{proof}
This is an autoformalization task, not a proof task. Produce statements that can later be proved by the physics prover without weakening the physical model.
\end{proof}

% --- Archon declaration topology begin ---
\section{Declaration topology}
\begin{definition}[Mass Quantity]
  \label{def:physics:phyx-mini-0704:phyxminiproblems-problemphyxmini0704-massquantity}
  \lean{PhyXMiniProblems.ProblemPhyXMini0704.MassQuantity}
  A nonnegative physical mass
\end{definition}
\begin{proof}
  This is the declared model definition of Mass Quantity.
\end{proof}
\begin{definition}[Length Quantity]
  \label{def:physics:phyx-mini-0704:phyxminiproblems-problemphyxmini0704-lengthquantity}
  \lean{PhyXMiniProblems.ProblemPhyXMini0704.LengthQuantity}
  A nonnegative physical length
\end{definition}
\begin{proof}
  This is the declared model definition of Length Quantity.
\end{proof}
\begin{definition}[Time Quantity]
  \label{def:physics:phyx-mini-0704:phyxminiproblems-problemphyxmini0704-timequantity}
  \lean{PhyXMiniProblems.ProblemPhyXMini0704.TimeQuantity}
  A nonnegative physical duration
\end{definition}
\begin{proof}
  This is the declared model definition of Time Quantity.
\end{proof}
\begin{definition}[Torque Magnitude Quantity]
  \label{def:physics:phyx-mini-0704:phyxminiproblems-problemphyxmini0704-torquemagnitudequantity}
  \lean{PhyXMiniProblems.ProblemPhyXMini0704.TorqueMagnitudeQuantity}
  A nonnegative torque magnitude, with dimension M L² T⁻²
\end{definition}
\begin{proof}
  This is the declared model definition of Torque Magnitude Quantity.
\end{proof}
\begin{definition}[Moment Of Inertia Quantity]
  \label{def:physics:phyx-mini-0704:phyxminiproblems-problemphyxmini0704-momentofinertiaquantity}
  \lean{PhyXMiniProblems.ProblemPhyXMini0704.MomentOfInertiaQuantity}
  A moment of inertia about the displayed axis, with dimension M L²
\end{definition}
\begin{proof}
  This is the declared model definition of Moment Of Inertia Quantity.
\end{proof}
\begin{definition}[Angular Speed Quantity]
  \label{def:physics:phyx-mini-0704:phyxminiproblems-problemphyxmini0704-angularspeedquantity}
  \lean{PhyXMiniProblems.ProblemPhyXMini0704.AngularSpeedQuantity}
  A nonnegative angular speed; radians are dimensionless
\end{definition}
\begin{proof}
  This is the declared model definition of Angular Speed Quantity.
\end{proof}
\begin{definition}[Angular Acceleration Quantity]
  \label{def:physics:phyx-mini-0704:phyxminiproblems-problemphyxmini0704-angularaccelerationquantity}
  \lean{PhyXMiniProblems.ProblemPhyXMini0704.AngularAccelerationQuantity}
  A nonnegative angular acceleration; radians are dimensionless
\end{definition}
\begin{proof}
  This is the declared model definition of Angular Acceleration Quantity.
\end{proof}
\begin{definition}[nonnegative SIReadout]
  \label{def:physics:phyx-mini-0704:phyxminiproblems-problemphyxmini0704-nonnegativesireadout}
  \lean{PhyXMiniProblems.ProblemPhyXMini0704.nonnegativeSIReadout}
  Read a nonnegative dimensionful quantity in coherent SI units
\end{definition}
\begin{proof}
  This is the declared model definition of nonnegative SIReadout.
\end{proof}
\begin{definition}[length In Meters]
  \label{def:physics:phyx-mini-0704:phyxminiproblems-problemphyxmini0704-lengthinmeters}
  \lean{PhyXMiniProblems.ProblemPhyXMini0704.lengthInMeters}
  \uses{def:physics:phyx-mini-0704:phyxminiproblems-problemphyxmini0704-lengthquantity, def:physics:phyx-mini-0704:phyxminiproblems-problemphyxmini0704-nonnegativesireadout}
  Metre readout of the total tip-to-tip propeller length
\end{definition}
\begin{proof}
  This is the declared model definition of length In Meters.
\end{proof}
\begin{definition}[mass In Kilograms]
  \label{def:physics:phyx-mini-0704:phyxminiproblems-problemphyxmini0704-massinkilograms}
  \lean{PhyXMiniProblems.ProblemPhyXMini0704.massInKilograms}
  \uses{def:physics:phyx-mini-0704:phyxminiproblems-problemphyxmini0704-massquantity, def:physics:phyx-mini-0704:phyxminiproblems-problemphyxmini0704-nonnegativesireadout}
  Kilogram readout of the propeller mass
\end{definition}
\begin{proof}
  This is the declared model definition of mass In Kilograms.
\end{proof}
\begin{definition}[time In Seconds]
  \label{def:physics:phyx-mini-0704:phyxminiproblems-problemphyxmini0704-timeinseconds}
  \lean{PhyXMiniProblems.ProblemPhyXMini0704.timeInSeconds}
  \uses{def:physics:phyx-mini-0704:phyxminiproblems-problemphyxmini0704-timequantity, def:physics:phyx-mini-0704:phyxminiproblems-problemphyxmini0704-nonnegativesireadout}
  Second readout of a physical duration
\end{definition}
\begin{proof}
  This is the declared model definition of time In Seconds.
\end{proof}
\begin{definition}[torque In Newton Meters]
  \label{def:physics:phyx-mini-0704:phyxminiproblems-problemphyxmini0704-torqueinnewtonmeters}
  \lean{PhyXMiniProblems.ProblemPhyXMini0704.torqueInNewtonMeters}
  \uses{def:physics:phyx-mini-0704:phyxminiproblems-problemphyxmini0704-torquemagnitudequantity, def:physics:phyx-mini-0704:phyxminiproblems-problemphyxmini0704-nonnegativesireadout}
  Newton-metre readout of a torque magnitude
\end{definition}
\begin{proof}
  This is the declared model definition of torque In Newton Meters.
\end{proof}
\begin{definition}[moment Of Inertia In Kilogram Meters Squared]
  \label{def:physics:phyx-mini-0704:phyxminiproblems-problemphyxmini0704-momentofinertiainkilogrammeterssquared}
  \lean{PhyXMiniProblems.ProblemPhyXMini0704.momentOfInertiaInKilogramMetersSquared}
  \uses{def:physics:phyx-mini-0704:phyxminiproblems-problemphyxmini0704-momentofinertiaquantity, def:physics:phyx-mini-0704:phyxminiproblems-problemphyxmini0704-nonnegativesireadout}
  Kilogram-metre-squared readout of a moment of inertia
\end{definition}
\begin{proof}
  This is the declared model definition of moment Of Inertia In Kilogram Meters Squared.
\end{proof}
\begin{definition}[angular Speed In Radians Per Second]
  \label{def:physics:phyx-mini-0704:phyxminiproblems-problemphyxmini0704-angularspeedinradianspersecond}
  \lean{PhyXMiniProblems.ProblemPhyXMini0704.angularSpeedInRadiansPerSecond}
  \uses{def:physics:phyx-mini-0704:phyxminiproblems-problemphyxmini0704-angularspeedquantity, def:physics:phyx-mini-0704:phyxminiproblems-problemphyxmini0704-nonnegativesireadout}
  Radian-per-second readout of an angular speed
\end{definition}
\begin{proof}
  This is the declared model definition of angular Speed In Radians Per Second.
\end{proof}
\begin{definition}[angular Acceleration In Radians Per Second Squared]
  \label{def:physics:phyx-mini-0704:phyxminiproblems-problemphyxmini0704-angularaccelerationinradianspersecondsquared}
  \lean{PhyXMiniProblems.ProblemPhyXMini0704.angularAccelerationInRadiansPerSecondSquared}
  \uses{def:physics:phyx-mini-0704:phyxminiproblems-problemphyxmini0704-angularaccelerationquantity, def:physics:phyx-mini-0704:phyxminiproblems-problemphyxmini0704-nonnegativesireadout}
  Radian-per-second-squared readout of an angular acceleration
\end{definition}
\begin{proof}
  This is the declared model definition of angular Acceleration In Radians Per Second Squared.
\end{proof}
\begin{definition}[angular Speed In Revolutions Per Minute]
  \label{def:physics:phyx-mini-0704:phyxminiproblems-problemphyxmini0704-angularspeedinrevolutionsperminute}
  \lean{PhyXMiniProblems.ProblemPhyXMini0704.angularSpeedInRevolutionsPerMinute}
  \uses{def:physics:phyx-mini-0704:phyxminiproblems-problemphyxmini0704-angularspeedquantity, def:physics:phyx-mini-0704:phyxminiproblems-problemphyxmini0704-angularspeedinradianspersecond}
  Revolutions-per-minute readout, using one revolution as 2 * π radians
\end{definition}
\begin{proof}
  This is the declared model definition of angular Speed In Revolutions Per Minute.
\end{proof}
\begin{definition}[Figure Object]
  \label{def:physics:phyx-mini-0704:phyxminiproblems-problemphyxmini0704-figureobject}
  \lean{PhyXMiniProblems.ProblemPhyXMini0704.FigureObject}
  Objects and graphical annotations visibly present in the supplied image
\end{definition}
\begin{proof}
  This is the declared model definition of Figure Object.
\end{proof}
\begin{definition}[Figure Text Label]
  \label{def:physics:phyx-mini-0704:phyxminiproblems-problemphyxmini0704-figuretextlabel}
  \lean{PhyXMiniProblems.ProblemPhyXMini0704.FigureTextLabel}
  Text labels printed in the supplied image
\end{definition}
\begin{proof}
  This is the declared model definition of Figure Text Label.
\end{proof}
\begin{definition}[Figure Point]
  \label{def:physics:phyx-mini-0704:phyxminiproblems-problemphyxmini0704-figurepoint}
  \lean{PhyXMiniProblems.ProblemPhyXMini0704.FigurePoint}
  Distinguished locations in the supplied propeller drawing
\end{definition}
\begin{proof}
  This is the declared model definition of Figure Point.
\end{proof}
\begin{definition}[Propeller Figure Readout]
  \label{def:physics:phyx-mini-0704:phyxminiproblems-problemphyxmini0704-propellerfigurereadout}
  \lean{PhyXMiniProblems.ProblemPhyXMini0704.PropellerFigureReadout}
  \uses{def:physics:phyx-mini-0704:phyxminiproblems-problemphyxmini0704-figureobject, def:physics:phyx-mini-0704:phyxminiproblems-problemphyxmini0704-figuretextlabel, def:physics:phyx-mini-0704:phyxminiproblems-problemphyxmini0704-figurepoint}
  Scalar labels and incidence information read directly from the image
\end{definition}
\begin{proof}
  This is the declared model definition of Propeller Figure Readout.
\end{proof}
\begin{definition}[Rotation Axis Placement]
  \label{def:physics:phyx-mini-0704:phyxminiproblems-problemphyxmini0704-rotationaxisplacement}
  \lean{PhyXMiniProblems.ProblemPhyXMini0704.RotationAxisPlacement}
  Placement of the fixed rotation axis relative to the propeller
\end{definition}
\begin{proof}
  This is the declared model definition of Rotation Axis Placement.
\end{proof}
\begin{definition}[Propeller Mass Model]
  \label{def:physics:phyx-mini-0704:phyxminiproblems-problemphyxmini0704-propellermassmodel}
  \lean{PhyXMiniProblems.ProblemPhyXMini0704.PropellerMassModel}
  Mass-distribution idealization used to compute the axial inertia
\end{definition}
\begin{proof}
  This is the declared model definition of Propeller Mass Model.
\end{proof}
\begin{definition}[Startup Torque Regime]
  \label{def:physics:phyx-mini-0704:phyxminiproblems-problemphyxmini0704-startuptorqueregime}
  \lean{PhyXMiniProblems.ProblemPhyXMini0704.StartupTorqueRegime}
  Torque regime during the modeled startup interval
\end{definition}
\begin{proof}
  This is the declared model definition of Startup Torque Regime.
\end{proof}
\begin{definition}[Propeller Startup Setup]
  \label{def:physics:phyx-mini-0704:phyxminiproblems-problemphyxmini0704-propellerstartupsetup}
  \lean{PhyXMiniProblems.ProblemPhyXMini0704.PropellerStartupSetup}
  \uses{def:physics:phyx-mini-0704:phyxminiproblems-problemphyxmini0704-massquantity, def:physics:phyx-mini-0704:phyxminiproblems-problemphyxmini0704-lengthquantity, def:physics:phyx-mini-0704:phyxminiproblems-problemphyxmini0704-timequantity, def:physics:phyx-mini-0704:phyxminiproblems-problemphyxmini0704-torquemagnitudequantity, def:physics:phyx-mini-0704:phyxminiproblems-problemphyxmini0704-momentofinertiaquantity, def:physics:phyx-mini-0704:phyxminiproblems-problemphyxmini0704-angularspeedquantity, def:physics:phyx-mini-0704:phyxminiproblems-problemphyxmini0704-angularaccelerationquantity, def:physics:phyx-mini-0704:phyxminiproblems-problemphyxmini0704-propellerfigurereadout, def:physics:phyx-mini-0704:phyxminiproblems-problemphyxmini0704-rotationaxisplacement, def:physics:phyx-mini-0704:phyxminiproblems-problemphyxmini0704-propellermassmodel, def:physics:phyx-mini-0704:phyxminiproblems-problemphyxmini0704-startuptorqueregime}
  The physical quantities and figure attached to the startup experiment
\end{definition}
\begin{proof}
  This is the declared model definition of Propeller Startup Setup.
\end{proof}
\begin{definition}[Matches Problem Data]
  \label{def:physics:phyx-mini-0704:phyxminiproblems-problemphyxmini0704-matchesproblemdata}
  \lean{PhyXMiniProblems.ProblemPhyXMini0704.MatchesProblemData}
  \uses{def:physics:phyx-mini-0704:phyxminiproblems-problemphyxmini0704-torqueinnewtonmeters, def:physics:phyx-mini-0704:phyxminiproblems-problemphyxmini0704-angularspeedinrevolutionsperminute, def:physics:phyx-mini-0704:phyxminiproblems-problemphyxmini0704-propellerstartupsetup}
  The values stated in the text for the driving torque and target speed
\end{definition}
\begin{proof}
  This is the declared model definition of Matches Problem Data.
\end{proof}
\begin{definition}[Matches Source Figure]
  \label{def:physics:phyx-mini-0704:phyxminiproblems-problemphyxmini0704-matchessourcefigure}
  \lean{PhyXMiniProblems.ProblemPhyXMini0704.MatchesSourceFigure}
  \uses{def:physics:phyx-mini-0704:phyxminiproblems-problemphyxmini0704-lengthinmeters, def:physics:phyx-mini-0704:phyxminiproblems-problemphyxmini0704-massinkilograms, def:physics:phyx-mini-0704:phyxminiproblems-problemphyxmini0704-propellerstartupsetup}
  The labels and geometry read from the supplied image, including the coherence of its scalar labels with the dimensionful setup quantities
\end{definition}
\begin{proof}
  This is the declared model definition of Matches Source Figure.
\end{proof}
\begin{definition}[Valid Propeller Startup Physics]
  \label{def:physics:phyx-mini-0704:phyxminiproblems-problemphyxmini0704-validpropellerstartupphysics}
  \lean{PhyXMiniProblems.ProblemPhyXMini0704.ValidPropellerStartupPhysics}
  \uses{def:physics:phyx-mini-0704:phyxminiproblems-problemphyxmini0704-lengthinmeters, def:physics:phyx-mini-0704:phyxminiproblems-problemphyxmini0704-massinkilograms, def:physics:phyx-mini-0704:phyxminiproblems-problemphyxmini0704-timeinseconds, def:physics:phyx-mini-0704:phyxminiproblems-problemphyxmini0704-torqueinnewtonmeters, def:physics:phyx-mini-0704:phyxminiproblems-problemphyxmini0704-momentofinertiainkilogrammeterssquared, def:physics:phyx-mini-0704:phyxminiproblems-problemphyxmini0704-angularspeedinradianspersecond, def:physics:phyx-mini-0704:phyxminiproblems-problemphyxmini0704-angularaccelerationinradianspersecondsquared, def:physics:phyx-mini-0704:phyxminiproblems-problemphyxmini0704-propellerstartupsetup}
  Governing fixed-axis dynamics and the uniform-rod idealization. The final startup duration is deliberately absent. The last field is the constant-angular-acceleration kinematic law relating the still-unknown duration to the initial and target angular speeds
\end{definition}
\begin{proof}
  This is the declared model definition of Valid Propeller Startup Physics.
\end{proof}
\begin{definition}[Answer Choice]
  \label{def:physics:phyx-mini-0704:phyxminiproblems-problemphyxmini0704-answerchoice}
  \lean{PhyXMiniProblems.ProblemPhyXMini0704.AnswerChoice}
  Letter labels used by the four displayed multiple-choice answers
\end{definition}
\begin{proof}
  This is the declared model definition of Answer Choice.
\end{proof}
\begin{definition}[duration Seconds]
  \label{def:physics:phyx-mini-0704:phyxminiproblems-problemphyxmini0704-answerchoice-durationseconds}
  \lean{PhyXMiniProblems.ProblemPhyXMini0704.AnswerChoice.durationSeconds}
  \uses{def:physics:phyx-mini-0704:phyxminiproblems-problemphyxmini0704-answerchoice}
  The duration in seconds printed beside each answer-choice letter
\end{definition}
\begin{proof}
  This is the declared model definition of duration Seconds.
\end{proof}
% --- Archon declaration topology end ---
% --- Archon physics formalization source end ---
